\input{\string~/Documents/School/"UC Irvine"/ta_template2.0.tex}
\rh{2B}{11.3}
\chead{\today}
\graphicspath{ {\string~/Desktop} }
\usetikzlibrary{arrows}
%============ANSWER KEY TOGGLE===========
\newcommand{\ans}[1]{ \noindent {\color{red} \textbf{Answer:} #1}}
\renewcommand{\ans}[1]{\vspace{3in}}
%==========================================
\begin{document}
\begin{center}
\textbf{The Integral Test}
\end{center}

\begin{thm}[Integral Test]
Suppose $f$ is a continuous, positive, decreasing function on $[1, \infty)$ and let $a_n=f(n)$. Then the series $\sum_{n=1}^{\infty} a_n$ is convergent if and only if the improper integral $\int_1^{\infty} f(x) d x$ is convergent. In other words:
\begin{enumerate}[label= \roman*.]
\item If $\int_1^{\infty} f(x) d x$ is convergent, then $\sum_{n=1}^{\infty} a_n$ is convergent.
\item If $\int_1^{\infty} f(x) d x$ is divergent, then $\sum_{n=1}^{\infty} a_n$ is divergent. 
\end{enumerate}
\end{thm}

\begin{aun}
For some intuition on the previous statement, consider drawing a Riemann sum of the integral with $\Delta x = 1$ and start computing the area beneath the curve that way. 
\end{aun}

\begin{thm}[$p$ Series Test]
The $p$-series 
\[
\sum_{n = 1}^\infty \frac{1}{n^p} \mte{ is } \begin{cases}
\text{convergent } & p > 1\\
\text{divergent } & p \leq 1
\end{cases}
\]
\end{thm}

\begin{exx}
Determine the convergence of the series:
\[
\sum_{n = 1}^\infty \frac{\ln(n)}{n}
\]
\end{exx}

\ans{In order to pair this up with the corresponding function: $f(x) = \frac{\ln(x)}{x}$, we have to establish all of the prerequisite conditions:
\begin{description}
\item[$f$ is Continuous:] $\ln(x), \frac1x$ are continuous on $(0,\infty)$, so their ratio is well defined and continuous. 
\item[$f$ is Positive:] Ratio of two positive functions on the domain $(0,\infty)$, so positive.
\item[$f$ is Decreasing:] Recall that it suffices to show that the derivative is negative. We see:
\[
f'(x) = \frac{\frac1x (x) - \ln(x) \cdot 1}{x^2} = \frac{1 - \ln(x)}{x^2} < 0 
\]
for all $x \in (0, \infty)$. 
\end{description}
These complete the conditions necessary for the Integral Test, so we can compute:
\[
\int_{1}^\infty \frac{\ln(x)}{x} dx
\]
Since this integral diverges, the series must diverge as well. 
}


\problembox{Compute the convergence of the following series:
\[
\sum_{n = 1}^\infty \frac{n}{n^2 + 1}
\]}

\ans{Corresponding integral diverges.}

\problembox{Compute the convergence of the following series:
\[
\sum_{n = 1}^\infty \frac{1}{n^2 + 2n + 2}
\]}

\ans{Converges, use the corresponding function, and then partial fractions to compute the integral. }

\problembox{Can we use the integral test on the following function?
\[
\sum_{n=1}^{\infty} \frac{\cos \pi n}{\sqrt{n}}
\]
}

\ans{No, positivity fails. }


\end{document}
